\mode*
\section{Method}

This study has two parts, mirroring the companion paper
\autocite{VTProgMisconceptions}.
The first is a \emph{literature review} that maps learners' misconceptions
and difficulties with Git, answering \cref{rq:misconceptions}:

\rqmisconceptions*

The second is an \emph{analysis} of those misconceptions through the lens
of variation theory: the misconceptions are treated as phenomenographic
observations from which we identify critical aspects (\cref{rq:aspects})
and construct patterns of variation (\cref{rq:patterns}):

\rqaspects*
\rqpatterns*

\Cref{rq:misconceptions} is answered by the review directly;
\cref{rq:aspects,rq:patterns} are answered analytically, and their answers
are validated by argument (inference to the best explanation), not yet
empirically.

\subsection{Literature review}
\label{sec:method-litreview}

The review has two phases.
The first phase is a seed round: a claim-driven search conducted for the
course-development work this paper grows out of, with every retained
reference verified against at least its abstract (full text where open)
and recorded with full provenance; the protocol and queries are in
\cref{app:protocol}.
The second phase is a systematic, tool-supported multi-database search
with the \texttt{scholar} tool, as in the companion papers
\autocite{VTProgMisconceptions,VTDebug}; it covered version-control
misconceptions beyond the seed sources, Git in CS-education venues (the
GitKit line and related), repository-data studies of student Git usage, and
Git visualisation tools as interventions.
Both rounds, with every query and what it retained, are documented in
\cref{app:protocol}.

\subsection{Variation-theory analysis}
\label{sec:method-vt-analysis}

For each documented misconception we ask: which aspect had the learner not
discerned?
Aspects recurring across misconceptions are candidate critical aspects.
For each critical aspect we construct patterns of variation---contrast,
generalisation, fusion---that make the aspect discernible, following the
companion paper \autocite{VTProgMisconceptions}.
Git affords the vary-and-observe loop that variation-theory teaching can
exploit \autocite{Svensson2020Affordances,ThuneEckerdal2009Variation,
ThuneEckerdal2018Lab}: the same command against different repository
states, or different commands against the same state, with
\texttt{git status}/\texttt{git log} as the observation instruments.

\subsection{Course instruments}
\label{sec:method-instruments}

The answers to \cref{rq:aspects,rq:patterns} are validated here by argument,
not yet empirically.
To enable a future empirical validation, we prepare a pre/post-test: a pair
of quizzes that measure whether students discern the candidate critical
aspects (\cref{sec:model,sec:staging,sec:branches,sec:remotes}) before and
after the course module that teaches against the misconceptions.
Each item targets one candidate aspect, and its distractors express the
documented misconceptions, so that a shift from a distractor to the key
across the pre- and post-test is evidence that the aspect became discernible.
The instrument is defined, item by item, as a literate program in
\cref{app:quiz}; running it and analysing the results is left to the
empirical follow-up.

The pre/post-test is one of two data sources the follow-up will have.
The other is the course's own assessment: the grading pipeline clones
every student's repository and runs per-repository checks, and each
check operationalises one of the candidate critical aspects---so each
check failure is an observation of an undiscerned aspect, collected by
the assessment the course already runs.
The pushed histories additionally carry process data (commit
granularity and timing, message content, branch use), the
repository-data perspective of \textcite{Cochez2013DVCSUsage} applied
to our own course; mining student repositories
\autocite{Hamer2021GitMetrics} and process-mining VCS interaction logs
\autocite{Ardimento2022ProcessMining} are established approaches we
build on.
The two sources are kept methodologically separate to avoid
circularity: misconception codes derived from the assessment traces are
never correlated with those same traces; the pre/post-test is the
independent measure of discernment.
Participation is opt-in through the consent item of the pre-test, which
covers course submissions; grading runs for all students regardless,
and research extraction filters by consent afterwards.
